\input{../tex-inputs/latex-leseansicht-vorspann}

\section[Arthur und Olga Schnitzler an Gustav Schwarzkopf, 4. 5. 1905]{L04386 Arthur und Olga Schnitzler an Gustav Schwarzkopf, 4. 5. 1905}
\nopagebreak\mylabel{L04386v}
\rehead{ }\normalsize\beginnumbering\briefempfaengerindex{Schwarzkopf, Gustav@\textsc{Schwarzkopf, Gustav}!zzzSchnitzler, Olga@\emph{von Olga Schnitzler}!1905-05-041@{4. 5. 1905}|(be}\briefempfaengerindex{Schwarzkopf, Gustav@\textsc{Schwarzkopf, Gustav}!zzzSchnitzler, Arthur@\emph{von Arthur Schnitzler}!1905-05-041@{4. 5. 1905}|(be}
\toendnotes[C]{\smallbreak\pagebreak[2]}
\correspDesc{Versand  durch Arthur Schnitzler, Olga Schnitzler am 4. 5. 1905 in Semmering
\newline{}Übermittlung  am 5. 5. 1905 in Semmering
\newline{}Zustellung  am 5. 5. 1905 in Wien 1/1 1
\newline{}Erhalt  durch Gustav Schwarzkopf am 5. 5. 1905 in Tiefer Graben 23}\toendnotes[C]{\smallbreak}
\Standort{DLA, A:Schnitzler, HS.1985.1.1897.}
\physDesc{Bildpostkarte, 142 Zeichen
\newline{}Handschrift Arthur Schnitzler: Bleistift, deutsche Kurrent
\newline{}Handschrift Olga Schnitzler: Bleistift
\newline{}Versand: 1) Stempel: »\nobreak{}\oindex{Semmering@\textbf{Semmering}|pwk}Semmering, 5. 5. 05, 6–7V\nobreak{}«.   2) Stempel: »\nobreak{}\oindex{I., Innere Stadt@\textbf{I., Innere Stadt}|pwk}Wien 1/1 1, 5 5 {[}05{]}, XI, Bestellt\nobreak{}«. }\pstart{}{\pb}Herrn \textsc{Gustav Schwarzkopf}\pend{}\pstart{}\textsc{Wien I}\oindex{I., Innere Stadt@\textbf{I., Innere Stadt}|pw}\pend{}\pstart{}\textsc{Tiefer Graben 23}\oindex{Wien@\textbf{Wien}!I., Innere Stadt@\textbf{I., Innere Stadt}!Tiefer Graben 23@\textbf{Tiefer Graben 23}|pw}.\pend{}{\bigskip}
\pstart
           \centering{}{\pb}\textcolor{gray}{\textbf{Semmering\oindex{Semmering@\textbf{Semmering}|pw}.}}\pend
           
\pstart
           \centering{}\textcolor{gray}{\textbf{Südbahnhôtel\oindex{Südbahnhotel [Semmering]@\textbf{Südbahnhotel [Semmering]}|pw}, 1000 m Seehöhe.}}\pend
           \vspace{1em}
\pstart
           {\pb}4. 5. 905.\pend
           \vspace{0.5em}
\pstart
           Herzlichen Gruß Ihnen und Bruder Max\pwindex{Schwarzkopf, Max 12.\,6.\,1857 Wien – 14.\,4.\,1928 ebd.@\textsc{Schwarzkopf, Max} (12.\,6.\,1857 Wien – 14.\,4.\,1928 ebd.)|pw} und auf
               baldgs Wiederſehen. Schön iſts hier!\pend
           \pstart \spacefill\mbox{A.}\pend{}\selectlanguage{ngerman}\vspace{1em}\pstart \spacefill\mbox{{[}hs. O. Schnitzler:{]} O. S.}\pend{}\selectlanguage{ngerman}\endnumbering\briefempfaengerindex{Schwarzkopf, Gustav@\textsc{Schwarzkopf, Gustav}!zzzSchnitzler, Olga@\emph{von Olga Schnitzler}!1905-05-041@{4. 5. 1905}|)be}\briefempfaengerindex{Schwarzkopf, Gustav@\textsc{Schwarzkopf, Gustav}!zzzSchnitzler, Arthur@\emph{von Arthur Schnitzler}!1905-05-041@{4. 5. 1905}|)be}\mylabel{L04386h}
\begin{anhang}
\end{anhang}\newcommand{\dateiname}{L04386}\newcommand{\titel}{Arthur und Olga Schnitzler an Gustav Schwarzkopf, 4. 5. 1905}\newcommand{\editorInnen}{Herausgegeben von Jahnke, SelmaMüller, Martin Anton}\input{../tex-inputs/latex-leseansicht-abspann}
